% How accurate LIGHT v0_3 is: v0_2 with the kinetic EEDF in the loop. The
% v0_2 overview describes the chemistry, which v0_3 keeps, and its rates
% (its Appendix A). No number here is typed by hand. They come from
%   overview_v0_3_fast.tex, overview_v0_3_slow.tex   make_overview_v0_3.py
%   overview_v0_2_fast.tex, overview_v0_2_slow.tex   make_overview_v0_2.py
%   energy_ladder_v0_3.tex                           make_energy_ladder.py
% Build from LIGHT/docs:
%   python3 make_overview_v0_3.py && python3 make_energy_ladder.py v0_3
%   latexmk -pdf accuracy_overview_v0_3.tex
\documentclass[11pt]{article}
\usepackage[a4paper,margin=2.2cm]{geometry}
\usepackage[T1]{fontenc}
\usepackage[scaled=0.95]{helvet}
\renewcommand{\familydefault}{\sfdefault}
\usepackage{amsmath}
\usepackage{sansmath}
\usepackage{booktabs,tabularx,array}
\usepackage{xcolor}
\usepackage{float}
\usepackage{enumitem}
\usepackage{needspace}
\usepackage{tikz}
\usetikzlibrary{arrows.meta,decorations.pathmorphing}
\setlist{nosep, leftmargin=1.2em}
\setlength{\parindent}{0pt}
\setlength{\parskip}{5pt}
\definecolor{ink2}{HTML}{52514E}

% LIGHT v0_2 accuracy numbers (fast). GENERATED by make_overview_v0_2.py
% from LIGHT v0_2 (0.1.2) and ML/. Do not edit. Rerun the script instead.
\providecommand\widef[2]{\expandafter\def\csname wi@#1\endcsname{#2}}
\widef{xu.50}{0.91}
\widef{xu.50.v0}{2.89}
\widef{xu.100}{0.97}
\widef{xu.100.v0}{2.41}
\widef{xu.50.meas}{$5.0{\times}10^{7}$}
\widef{xu.100.meas}{$6.2{\times}10^{7}$}
\widef{eq5.dev}{$1.7{\times}10^{-4}$}
\widef{gap}{120}
\widef{closure.max}{$10^{-12}$}
\widef{closure.runs}{9}
\widef{hot.p}{0.1}
\widef{hot.P}{500}
\widef{hot.tg}{3458}
\widef{sc.eden_kr}{2.3}
\widef{sc.eden_f2}{3.2}
\widef{sc.ksu_fig2}{4.2}
\widef{sc.ksu_fig1}{6.1}
\widef{sc.xu_figs5_7}{11.7}
\widef{sc.hakuta_fig3}{34.1}
\widef{sc.hakuta_fig2}{19.3}
\widef{sc.ys_ir300}{\ensuremath{-}0.5}
\widef{sc.ys_ratio_230_300}{\ensuremath{-}8.4}
\widef{sc.bs_tau_b}{\ensuremath{-}22.3}
\widef{ys.meas}{0.16}
\widef{ys.err}{0.02}
\widef{ysr.meas}{1.30}
\widef{ysr.err}{0.06}
\widef{sc.train.lo}{2.3}
\widef{sc.train.hi}{6.1}
\widef{sc.bs.meas}{9.0}
\widef{sc.bs.err}{0.6}
\widef{tah.g0.meas}{6.39}
\widef{tah.g0}{5.69}
\widef{tah.g0.x}{0.89}
\widef{tah.isat.meas}{1.4}
\widef{tah.isat}{1.55}
\widef{tah.isat.err}{0.1}
\widef{tah.isat.x}{1.11}
\widef{tah.eta.meas}{0.27}
\widef{tah.eta}{0.268}
\widef{tah.eta.x}{0.99}
\widef{tah.form}{0.300}
\widef{tah.tg}{446}
\widef{tah.pin}{330}
\widef{tah.torr}{665}
\widef{tah.f2torr}{2.8}
\widef{tah.seed}{$1.0{\times}10^{10}$}
\widef{tah.def.g0}{5.69}
\widef{tah.win.lo}{200}
\widef{tah.win.hi}{400}
\widef{tah.f2left}{29}
\widef{tah.bd.kr2_f.lo}{5.69}
\widef{tah.bd.kr2_f.hi}{6.23}
\widef{tah.bd.trimer_tg_pow.lo}{5.69}
\widef{tah.bd.trimer_tg_pow.hi}{2.84}
\widef{tah.bd.kr_pen.lo}{5.71}
\widef{tah.bd.kr_pen.hi}{5.69}
\widef{tah.bd.kr2f_f2.lo}{5.69}
\widef{tah.bd.kr2f_f2.hi}{5.89}
\widef{tah.bd.kr2f_f.lo}{5.69}
\widef{tah.bd.kr2f_f.hi}{5.86}
\widef{tah.bd.kr2_form.lo}{5.69}
\widef{tah.bd.kr2_form.hi}{5.31}
\widef{tah.bd.kr2_f2.lo}{5.69}
\widef{tah.bd.kr2_f2.hi}{5.74}
\widef{tah.bd.kr2_pen.lo}{5.69}
\widef{tah.bd.kr2_pen.hi}{5.69}
\widef{tah.bd.minor}{0.38}
\widef{tah.bd.widest}{the trimer's $(T_g/300)^n$}
\widef{tah.hiF.g0}{6.23}
\widef{tah.hiF.x}{0.97}
\widef{tah.tri.g0}{2.84}
\widef{tah.tri.x}{0.44}
\widef{tri.lo.pow}{0.5}
\widef{tri.klimek.hi}{425}
\widef{tah.kannari.g0}{6.10}
\widef{tah.kannari.x}{0.95}
\widef{tah.own.g0}{5.53}
\widef{tah.own.te}{3.8}
\widef{tah.own.ne}{$2.59{\times}10^{14}$}
\widef{tah.own.tg}{423}
\widef{tah.beam.g0}{5.69}
\widef{tah.beam.te}{1.6}
\widef{tah.beam.ne}{$6.33{\times}10^{14}$}
\widef{tah.beam.tg}{446}
\widef{tah.own.x}{0.87}
\widef{kr2.def}{4.67}
\widef{kr2.vuv}{20.0}
\widef{kr2.burn}{640}
\widef{kr2.F}{$2.12{\times}10^{17}$}
\widef{kr2.lo.k}{$1.0{\times}10^{-11}$}
\widef{kr2.lo}{6.52}
\widef{kr2.kann.k}{$3.0{\times}10^{-10}$}
\widef{kr2.kann}{9.21}
\widef{kr2.kann.ratio}{17}
\widef{kr2.hi.k}{$1.0{\times}10^{-9}$}
\widef{kr2.hi}{9.50}
\widef{kr2.hi.vuv}{0.5}
\widef{kr2.hi.ratio}{57}
\widef{kr2.hi.afterpc}{74}
\widef{kr2.both}{9.17}
\widef{kr2.pen.k}{$5.0{\times}10^{-10}$}
\widef{pure.Kr.x}{1.02}
\widef{pure.Ar.x}{1.03}
\widef{pure.dr.lo}{1.00}
\widef{pure.dr.hi}{1.00}
\widef{pure.dr.te.lo}{0.026}
\widef{pure.dr.te.hi}{2}
\widef{pure.k15}{$2.4{\times}10^{-31}$}
\widef{pure.conv.lo}{$2.74{\times}10^{-31}$}
\widef{pure.conv.hi}{$3.3{\times}10^{-31}$}
\widef{pure.vuv}{0.57}
\widef{pure.vuv.meas}{0.42}
\widef{pure.vuv.err}{0.05}
\widef{pure.vuv.x}{1.36}
\widef{pure.vuv.ideal}{0.54}
\widef{pure.vuv.nothermal}{0.55}
\widef{pure.thermal}{0.060}
\widef{pure.band.move}{0.001}
\widef{pure.hi.P}{300}
\widef{pure.hi.pulse}{1}
\widef{pure.hi.torr}{665}
\widef{pure.hi.def}{0.27}
\widef{pure.hi.lo}{0.13}
\widef{pure.hi.lo.what}{\Krts{} + \Krts{} \ra \Krtp{} + 2Kr + \el{} at $5.0{\times}10^{-10}$\unit{cm^3\,s^{-1}}}
\widef{pure.hi.hi}{0.49}
\widef{pure.hi.hi.what}{\Krs{} + 2Kr \ra \Krts{} + Kr at $5.5{\times}10^{-32}$\unit{cm^6\,s^{-1}}}
\def\wiTahBands{\Krts{} + F \ra KrF(B) + Kr & 5.69 to 6.23 & \x{0.89} to \x{0.97}\\
the trimer's $(T_g/300)^n$ & 5.69 to 2.84 & \x{0.89} to \x{0.44}\\
\Krs{} + \Krs{} \ra \Krp{} + Kr + \el{} & 5.71 to 5.69 & \x{0.89} to \x{0.89}\\
\KtF{} + \Ft{} & 5.69 to 5.89 & \x{0.89} to \x{0.92}\\
\KtF{} + F \ra 2Kr + 2F & 5.69 to 5.86 & \x{0.89} to \x{0.92}\\
\Krs{} + 2Kr \ra \Krts{} + Kr, 300\unit{K} & 5.69 to 5.31 & \x{0.89} to \x{0.83}\\
\Krts{} + \Ft{}, both channels & 5.69 to 5.74 & \x{0.89} to \x{0.90}\\
\Krts{} + \Krts{} \ra \Krtp{} + 2Kr + \el{} & 5.69 to 5.69 & \x{0.89} to \x{0.89}\\}

% LIGHT v0_2 accuracy numbers (slow). GENERATED by make_overview_v0_2.py
% from LIGHT v0_2 (0.1.2) and ML/. Do not edit. Rerun the script instead.
\providecommand\widef[2]{\expandafter\def\csname wi@#1\endcsname{#2}}
\widef{fit.k_BC.lo}{$4.63{\times}10^{-12}$}
\widef{fit.k_BC.mid}{$9.8{\times}10^{-12}$}
\widef{fit.k_BC.hi}{$1.73{\times}10^{-11}$}
\widef{fit.K_eq.lo}{0.84}
\widef{fit.K_eq.mid}{1.07}
\widef{fit.K_eq.hi}{1.38}
\widef{fit.tau_C.lo}{41}
\widef{fit.tau_C.mid}{50}
\widef{fit.tau_C.hi}{63}
\widef{fit.tau_B.lo}{6.79}
\widef{fit.tau_B.mid}{7.02}
\widef{fit.tau_B.hi}{7.22}
\widef{fit.k_F2.lo}{$3.92{\times}10^{-10}$}
\widef{fit.k_F2.mid}{$4.57{\times}10^{-10}$}
\widef{fit.k_F2.hi}{$5.3{\times}10^{-10}$}
\widef{fit.k_F2_C.lo}{$2.74{\times}10^{-10}$}
\widef{fit.k_F2_C.mid}{$3.02{\times}10^{-10}$}
\widef{fit.k_F2_C.hi}{$3.32{\times}10^{-10}$}
\widef{fit.k3_Kr.lo}{$1.26{\times}10^{-30}$}
\widef{fit.k3_Kr.mid}{$1.92{\times}10^{-30}$}
\widef{fit.k3_Kr.hi}{$2.44{\times}10^{-30}$}
\widef{fit.k3_Kr_C.lo}{$1.93{\times}10^{-31}$}
\widef{fit.k3_Kr_C.mid}{$3.64{\times}10^{-31}$}
\widef{fit.k3_Kr_C.hi}{$5.12{\times}10^{-31}$}
\widef{fit.bshare.lo}{0.57}
\widef{fit.bshare.mid}{0.69}
\widef{fit.bshare.hi}{0.80}
\widef{fit.bshare.light}{0.68}
\widef{fit.n}{10}
\widef{post.g0.lo}{4.99}
\widef{post.g0.mid}{5.39}
\widef{post.g0.hi}{6.18}
\widef{post.n}{40}
\widef{link.base}{0.268}
\widef{link.yield.lo}{0.240}
\widef{link.yield.lo.d}{\ensuremath{-}0.028}
\widef{link.yield.hi}{0.282}
\widef{link.yield.hi.d}{+0.014}
\widef{link.yield.arf}{0.211}
\widef{link.yield.arf.d}{\ensuremath{-}0.057}
\widef{link.nascentC}{0.290}
\widef{link.nascentC.d}{+0.022}
\widef{link.ionY}{0.225}
\widef{link.ionY.d}{\ensuremath{-}0.043}
\widef{link.ionC}{0.251}
\widef{link.ionC.d}{\ensuremath{-}0.017}
\widef{link.ionRate}{0.266}
\widef{link.ionRate.d}{\ensuremath{-}0.002}
\widef{link.kr2F.kann}{0.285}
\widef{link.kr2F.kann.d}{+0.017}
\widef{link.kr2F.hi}{0.290}
\widef{link.kr2F.hi.d}{+0.022}
\widef{link.kr2split.lo}{0.260}
\widef{link.kr2split.lo.d}{\ensuremath{-}0.008}
\widef{link.kr2split.hi}{0.276}
\widef{link.kr2split.hi.d}{+0.008}
\widef{link.k3}{0.256}
\widef{link.k3.d}{\ensuremath{-}0.012}
\widef{link.minor}{0.013}
\widef{el.loop.te}{1.61}
\widef{el.loop.ne}{$6.49{\times}10^{14}$}
\widef{el.loop.g0}{5.63}
\widef{el.loop.form}{0.299}
\widef{el.light.te}{1.60}
\widef{el.light.ne}{$6.33{\times}10^{14}$}
\widef{el.light.g0}{5.69}
\widef{el.light.form}{0.300}
\widef{el.dev.te}{0.7}
\widef{el.dev.ne}{2.5}
\widef{el.dev.g0}{1.0}
\widef{el.dev.formation}{0.5}
\widef{el.eatt}{0.55}
\widef{el.eatt.32}{2.41}
\widef{el.rpow}{0.55}
\widef{el.rcap}{0.95}
\widef{sw.krw.lo}{0.96}
\widef{sw.krw.hi}{1.03}
\widef{sw.f2.lo}{0.48}
\widef{sw.f2.hi}{2.41}
\def\wiLinks{KrF$^*$ yield of \Krs{} + \Ft & 0.90 & 0.70 & the measured 0.90 less its stated 0.20 & 0.240 & \ensuremath{-}0.028\\
 &  & 1.0 & the published assignment & 0.282 & +0.014\\
 &  & 0.5 & as for Ar$^*$ + F$_2$ & 0.211 & \ensuremath{-}0.057\\
share of that KrF formed in C & 0.38 & 0 & if nascent C relaxes to B at 1 atm & 0.290 & +0.022\\
KrF$^*$ yield of ion-ion recombination & 1 & 0.64 & as for Cl$^-$ & 0.225 & \ensuremath{-}0.043\\
share of that KrF formed in C & 0 & 0.38 & as for SF$_6^-$ & 0.251 & \ensuremath{-}0.017\\
ion-ion recombination rate & Flannery & \x{0.38} & Avco's model input & 0.266 & \ensuremath{-}0.002\\
\Krts{} + F \ra KrF(B) + Kr & 0 & $3.0{\times}10^{-10}$ & Kannari's estimate & 0.285 & +0.017\\
 &  & $1.0{\times}10^{-9}$ & the top of its band & 0.290 & +0.022\\}

% LIGHT v0_3 accuracy numbers (fast). GENERATED by make_overview_v0_3.py
% from LIGHT v0_3 (0.1.3) and ML/. Do not edit. Rerun the script instead.
\providecommand\wkdef[2]{\expandafter\def\csname wk@#1\endcsname{#2}}
\wkdef{sc.eden_kr}{2.3}
\wkdef{sc.eden_f2}{3.2}
\wkdef{sc.ksu_fig2}{4.2}
\wkdef{sc.ksu_fig1}{6.1}
\wkdef{sc.xu_figs5_7}{11.7}
\wkdef{sc.hakuta_fig3}{34.1}
\wkdef{sc.hakuta_fig2}{19.3}
\wkdef{sc.ys_ir300}{\ensuremath{-}0.5}
\wkdef{sc.ys_ratio_230_300}{\ensuremath{-}8.4}
\wkdef{sc.bs_tau_b}{\ensuremath{-}22.3}
\wkdef{sc.same}{10}
\wkdef{sc.train.lo}{2.3}
\wkdef{sc.train.hi}{6.1}
\wkdef{xu.50}{0.91}
\wkdef{xu.100}{0.97}
\wkdef{ml.dF}{$2.2{\times}10^{-9}$}
\wkdef{ml.dk}{$4.2{\times}10^{-8}$}
\wkdef{ml.own}{$1.1{\times}10^{-5}$}
\wkdef{ml.states}{3}
\wkdef{bal.pulse}{$2.0{\times}10^{-6}$}
\wkdef{bal.pulse.nu}{0}
\wkdef{bal.nofield}{$5.6{\times}10^{-5}$}
\wkdef{bal.nofield.nu}{0}
\wkdef{bal.nogrowth}{$2.9{\times}10^{-4}$}
\wkdef{bal.nogrowth.raw}{$2.0{\times}10^{-3}$}
\wkdef{und.600.dev}{$1.4{\times}10^{-14}$}
\wkdef{und.600.cell}{0.029}
\wkdef{und.600.mean}{0.0377}
\wkdef{und.2400.dev}{$6.3{\times}10^{-14}$}
\wkdef{und.2400.cell}{0.007}
\wkdef{und.2400.mean}{0.0386}
\wkdef{und.32kt}{0.0388}
\wkdef{s.hot}{0.282}
\wkdef{s.cold}{0.042}
\wkdef{s.en.006}{0.444}
\wkdef{s.en.02}{0.388}
\wkdef{s.en.05}{0.473}
\wkdef{inj.dev}{$1.6{\times}10^{-5}$}
\wkdef{step.big}{$1.7{\times}10^{-12}$}
\wkdef{step.tiny}{$1.2{\times}10^{-10}$}
\wkdef{step.energy}{$2.8{\times}10^{-5}$}
\wkdef{step.ee600}{$1.0{\times}10^{-1}$}
\wkdef{step.ee2400}{$7.6{\times}10^{-3}$}
\wkdef{step.power}{$10^{9}$}
\wkdef{ref.p}{1}
\wkdef{ref.xf2}{0.5}
\wkdef{ref.P}{300}
\wkdef{ref.pulse}{1}
\wkdef{ref.tend}{2}
\wkdef{conv.dt.ns}{5}
\wkdef{und.lo}{600}
\wkdef{und.hi}{2400}
\wkdef{und.tg}{300}
\wkdef{s.t}{0.12}
\wkdef{s.field}{0.4}
\wkdef{s.m.006}{0.06}
\wkdef{s.m.02}{0.2}
\wkdef{s.m.05}{0.5}

% LIGHT v0_3 accuracy numbers (slow). GENERATED by make_overview_v0_3.py
% from LIGHT v0_3 (0.1.3) and ML/. Do not edit. Rerun the script instead.
\providecommand\wkdef[2]{\expandafter\def\csname wk@#1\endcsname{#2}}
\wkdef{ref.248}{4.79}
\wkdef{ref.248.v02}{4.67}
\wkdef{ref.248.d}{+2.7}
\wkdef{ref.272}{0.153}
\wkdef{ref.272.v02}{0.149}
\wkdef{ref.272.d}{+2.7}
\wkdef{ref.275}{0.735}
\wkdef{ref.275.v02}{0.712}
\wkdef{ref.275.d}{+3.2}
\wkdef{ref.420}{4.96}
\wkdef{ref.420.v02}{4.87}
\wkdef{ref.420.d}{+1.8}
\wkdef{ref.147}{19.7}
\wkdef{ref.147.v02}{20.0}
\wkdef{ref.147.d}{\ensuremath{-}1.8}
\wkdef{ref.all}{30.3}
\wkdef{ref.all.v02}{30.4}
\wkdef{ref.tg}{650}
\wkdef{ref.tg.v02}{649}
\wkdef{ref.krstar}{99.4}
\wkdef{ref.krstar.v02}{98.3}
\wkdef{ref.krplus}{26.0}
\wkdef{ref.krplus.v02}{26.8}
\wkdef{ref.krfb}{24.6}
\wkdef{ref.krfb.v02}{23.8}
\wkdef{ref.elastic}{5.3}
\wkdef{ref.elastic.v02}{6.2}
\wkdef{ref.closure}{$10^{-13}$}
\wkdef{ref.solutions}{355}
\wkdef{ref.en.lo}{0.98}
\wkdef{ref.en.hi}{43}
\wkdef{ref.ratio.lo}{1.006}
\wkdef{ref.ratio.hi}{1.031}
\wkdef{ref.gap}{$2.9{\times}10^{-5}$}
\wkdef{conv.step}{$3.9{\times}10^{-5}$}
\wkdef{conv.dt}{$6.8{\times}10^{-5}$}
\wkdef{conv.cells}{$1.6{\times}10^{-4}$}
\wkdef{ref.ratio.dt.hi}{1.015}
\wkdef{first.d.steady}{$1.9{\times}10^{-4}$}
\wkdef{first.d.maxwell}{$9.6{\times}10^{-4}$}
\wkdef{first.n}{87}
\wkdef{first.steady.n}{88}
\wkdef{first.steady.goff}{86}
\wkdef{first.mem.goff}{0}
\wkdef{first.mem.unreached}{0}
\wkdef{first.steady.ratio.lo}{0.02}
\wkdef{first.steady.ratio.hi}{0.19}
\wkdef{first.mem.ratio.lo}{0.98}
\wkdef{first.mem.ratio.hi}{1.07}
\wkdef{first.steady.gap}{0.22}
\wkdef{first.mem.gap0}{0.07}
\wkdef{first.mem.gap}{$1.5{\times}10^{-3}$}
\wkdef{first.edep}{$1.0{\times}10^{-3}$}
\wkdef{first.t}{1}
\wkdef{first.t.split}{0.1}
\wkdef{closure.max}{$10^{-13}$}
\wkdef{closure.runs}{4}
\wkdef{fac.ref.te}{3.00}
\wkdef{fac.ref.t}{502}
\wkdef{fac.tah.te}{1.61}
\wkdef{fac.tah.t}{302}
\wkdef{tah.g0}{5.64}
\wkdef{tah.g0.x}{0.88}
\wkdef{tah.g0.meas}{6.39}
\wkdef{tah.isat}{1.56}
\wkdef{tah.isat.x}{1.11}
\wkdef{tah.isat.meas}{1.4}
\wkdef{tah.eta}{0.266}
\wkdef{tah.eta.x}{0.99}
\wkdef{tah.eta.meas}{0.27}
\wkdef{tah.form}{0.299}
\wkdef{tah.win.g0}{5.64}
\wkdef{tah.win.te}{1.61}
\wkdef{tah.win.ne}{$6.52{\times}10^{14}$}
\wkdef{tah.win.tg}{446}
\wkdef{tah.gap}{$6.0{\times}10^{-4}$}
\wkdef{tah.reached}{every}
\wkdef{tah.en}{0.13}
\wkdef{first.tah.d}{$3.6{\times}10^{-4}$}
\wkdef{first.tah.ne.d}{$9.2{\times}10^{-4}$}
\wkdef{first.tah.steady.unreached}{80}
\wkdef{first.tah.mem.unreached}{0}
\wkdef{tah.kannari.g0}{6.05}
\wkdef{tah.kannari.x}{0.95}
\wkdef{tah.bd.kr2_f.lo}{5.64}
\wkdef{tah.bd.kr2_f.hi}{6.18}
\wkdef{tah.bd.trimer_tg_pow.lo}{5.64}
\wkdef{tah.bd.trimer_tg_pow.hi}{2.82}
\wkdef{tah.bd.kr_pen.lo}{5.66}
\wkdef{tah.bd.kr_pen.hi}{5.64}
\wkdef{tah.bd.kr2f_f2.lo}{5.64}
\wkdef{tah.bd.kr2f_f2.hi}{5.84}
\wkdef{tah.bd.kr2f_f.lo}{5.64}
\wkdef{tah.bd.kr2f_f.hi}{5.81}
\wkdef{tah.bd.kr2_form.lo}{5.64}
\wkdef{tah.bd.kr2_form.hi}{5.26}
\wkdef{tah.bd.kr2_f2.lo}{5.64}
\wkdef{tah.bd.kr2_f2.hi}{5.69}
\wkdef{tah.bd.kr2_pen.lo}{5.64}
\wkdef{tah.bd.kr2_pen.hi}{5.64}
\wkdef{tah.bd.widest}{the trimer's $(T_g/300)^n$}
\wkdef{tah.hiF.g0}{6.18}
\wkdef{tah.hiF.x}{0.97}
\wkdef{tah.tri.g0}{2.82}
\wkdef{tah.tri.x}{0.44}
\wkdef{el.loop.te}{1.610}
\wkdef{el.loop.ne}{$6.49{\times}10^{14}$}
\wkdef{el.loop.g0}{5.632}
\wkdef{el.loop.form}{0.2988}
\wkdef{el.own.te}{1.609}
\wkdef{el.own.ne}{$6.52{\times}10^{14}$}
\wkdef{el.own.g0}{5.637}
\wkdef{el.own.form}{0.2990}
\wkdef{el.dev.te}{\ensuremath{-}0.08}
\wkdef{el.abs.te}{0.08}
\wkdef{el.dev.ne}{+0.47}
\wkdef{el.abs.ne}{0.47}
\wkdef{el.dev.g0}{+0.08}
\wkdef{el.abs.g0}{0.08}
\wkdef{el.dev.formation}{+0.07}
\wkdef{el.abs.formation}{0.07}
\wkdef{el.passes}{10}
\wkdef{pure.vuv}{0.55}
\wkdef{pure.vuv.meas}{0.42}
\wkdef{pure.vuv.err}{0.05}
\wkdef{pure.vuv.x}{1.32}
\wkdef{pure.hi}{0.28}
\wkdef{pure.hi.P}{300}
\def\wkFactors{G1 & 0.076 & 0.070\\
C1a & 0.013 & $4.8{\times}10^{-3}$\\
C1b & $3.4{\times}10^{-3}$ & $7.0{\times}10^{-4}$\\
C1c & $1.5{\times}10^{-4}$ & $1.2{\times}10^{-5}$\\
C2 & $2.0{\times}10^{-5}$ & $2.3{\times}10^{-7}$\\
ATT & 0.58 & 0.96\\
DIS\_a3 & 1.2 & 1.0\\
DIS\_A1 & 1.1 & 0.97\\
VIB1 & 0.65 & 0.97\\
VIB4 & 0.69 & 0.98\\
R11p & 0.87 & 0.98\\}
\def\wkTahBands{\Krts{} + F \ra KrF(B) + Kr & 5.64 to 6.18 & \x{0.88} to \x{0.97}\\
the trimer's $(T_g/300)^n$ & 5.64 to 2.82 & \x{0.88} to \x{0.44}\\
\Krs{} + \Krs{} \ra \Krp{} + Kr + \el{} & 5.66 to 5.64 & \x{0.89} to \x{0.88}\\
\KtF{} + \Ft{} & 5.64 to 5.84 & \x{0.88} to \x{0.91}\\
\KtF{} + F \ra 2Kr + 2F & 5.64 to 5.81 & \x{0.88} to \x{0.91}\\
\Krs{} + 2Kr \ra \Krts{} + Kr, 300\unit{K} & 5.64 to 5.26 & \x{0.88} to \x{0.82}\\
\Krts{} + \Ft{}, both channels & 5.64 to 5.69 & \x{0.88} to \x{0.89}\\
\Krts{} + \Krts{} \ra \Krtp{} + 2Kr + \el{} & 5.64 to 5.64 & \x{0.88} to \x{0.88}\\}

% LIGHT v0_3 energy ladder. GENERATED by make_energy_ladder.py
% from LIGHT v0_3 (0.1.3). Do not edit. Rerun the script instead.
\providecommand\laddef[2]{\expandafter\def\csname lad@#1\endcsname{#2}}
\laddef{version}{0.1.3}
\laddef{E.B}{5.856}
\laddef{DeB}{5.54}
\laddef{BoverKrF}{5.06}
\laddef{wl.G4}{248}
\laddef{hv.G4}{5.00}
\laddef{heat.G4}{0.06}
\laddef{heat.G3}{3.26}
\laddef{heat.K26}{5.54}
\laddef{heat.G5}{0.58}
\laddef{wl.C3}{420}
\laddef{hv.C3}{2.95}
\laddef{heat.C3}{1.53}
\laddef{gapBC}{120}
\laddef{notdrawn}{G6, VIB1, VIB2, VIB3 and VIB4}
\laddef{toe}{G2, PEN\_KR2 and SUPER}
\laddef{zero}{Q2FF, K2F\_F and PEN\_KR2}
\laddef{groups}{\texttt{electron}}
\laddef{pale}{}
\laddef{wl.V}{147}
\laddef{hv.V}{8.43}
\laddef{heat.V}{0.77}
\def\ladpicture{%
\definecolor{ladE}{HTML}{2A78D6}% electron gas
\definecolor{ladX}{HTML}{EDA100}% Kr*
\definecolor{ladB}{HTML}{1BAF7A}% KrF(B), 248 and 272 nm
\definecolor{ladC}{HTML}{008300}% KrF(C), 275 nm
\definecolor{ladT}{HTML}{4A3AA7}% Kr2F*, 420 nm
\definecolor{ladV}{HTML}{A63A83}% Kr2*, 147 nm
\definecolor{ladH}{HTML}{E34948}% heat
\definecolor{ladI}{HTML}{52514E}% ions
\definecolor{ladF}{HTML}{A8A69E}% F atoms
\definecolor{ladInk}{HTML}{0B0B0B}%
\definecolor{ladInk2}{HTML}{52514E}%
\definecolor{ladMuted}{HTML}{898781}%
\tikzset{
  lad grid/.style={draw=ladMuted!22, line width=0.3pt},
  lad axisline/.style={draw=ladMuted, line width=0.5pt},
  lad lev/.style={line width=1.4pt},
  lad pair/.style={line width=0.8pt, dash pattern=on 3pt off 2pt},
  lad tick/.style={line width=0.5pt},
  lad arr/.style={line width=0.7pt, -{Triangle[width=3.2pt,length=3.4pt]}},
  lad hv/.style={line width=1.3pt, -{Triangle[width=4.2pt,length=4.4pt]}},
  lad on/.style={},
  lad opt/.style={dash pattern=on 2.4pt off 1.6pt},
  lad zero/.style={dash pattern=on 0.6pt off 1.4pt, line cap=round},
  lad heat/.style={draw=ladH, line width=0.6pt, decorate,
    decoration={snake, amplitude=0.5mm, segment length=1.4mm,
    post length=0.5mm}},
  lad lab/.style={font=\scriptsize, text=ladInk, inner sep=1pt},
  lad lab2/.style={font=\scriptsize, text=ladInk, inner sep=1pt,
    align=right},
  lad lab2l/.style={font=\scriptsize, text=ladInk, inner sep=1pt,
    align=left},
  lad alab/.style={font=\scriptsize, text=ladInk, inner sep=0.8pt,
    fill=white, fill opacity=0.85, text opacity=1, align=center},
  lad slab/.style={font=\tiny, text=ladInk, inner sep=0.6pt,
    fill=white, fill opacity=0.85, text opacity=1, align=right},
  lad hlab/.style={font=\tiny, text=ladInk2, inner sep=0.8pt,
    fill=white, fill opacity=0.85, text opacity=1},
  lad tlab/.style={font=\tiny, text=ladInk2, inner sep=0.8pt},
  lad axis/.style={font=\scriptsize, text=ladInk2, inner sep=1pt},
  lad head/.style={font=\footnotesize\bfseries, text=ladInk2},
}
\begin{tikzpicture}
\draw[lad grid] (0.950,1.000) -- (16.200,1.000);
\draw[lad grid] (0.950,3.000) -- (16.200,3.000);
\draw[lad grid] (0.950,5.000) -- (16.200,5.000);
\draw[lad grid] (0.950,7.000) -- (16.200,7.000);
\draw[lad grid] (0.950,9.000) -- (16.200,9.000);
\draw[lad grid] (0.950,11.000) -- (16.200,11.000);
\draw[lad grid] (0.950,13.000) -- (16.200,13.000);
\draw[lad grid] (0.950,15.000) -- (16.200,15.000);
\draw[lad grid] (0.950,17.000) -- (16.200,17.000);
\draw[lad axisline] (0.950,0.000) -- (0.950,17.450);
\draw[lad lev,draw=ladInk2] (2.650,3.000) -- (6.300,3.000);
\draw[lad lev,draw=ladX] (2.650,12.915) -- (6.300,12.915);
\draw[lad lev,draw=ladI] (2.650,16.996) -- (6.300,16.996);
\draw[lad lev,draw=ladV] (5.000,12.200) -- (6.300,12.200);
\draw[lad lev,draw=ladI] (5.000,15.846) -- (6.300,15.846);
\draw[lad tick,draw=ladMuted] (3.550,14.300) -- (4.400,14.300);
\draw[lad tick,draw=ladMuted] (3.730,14.998) -- (4.400,14.998);
\draw[lad tick,draw=ladMuted] (3.910,15.750) -- (4.400,15.750);
\draw[lad tick,draw=ladMuted] (5.740,3.770) -- (6.020,3.770);
\draw[lad pair,draw=ladInk2] (10.200,14.396) -- (12.300,14.396);
\draw[lad pair,draw=ladInk2] (10.200,13.246) -- (12.300,13.246);
\draw[lad pair,draw=ladInk2] (7.000,3.800) -- (12.300,3.800);
\draw[lad lev,draw=ladB] (8.400,8.886) -- (11.050,8.886);
\draw[lad lev,draw=ladC] (8.400,8.811) -- (11.050,8.811);
\draw[lad lev,draw=ladT] (7.300,8.276) -- (8.650,8.276);
\draw[lad tick,draw=ladMuted] (9.210,3.856) -- (9.490,3.856);
\draw[lad tick,draw=ladMuted] (9.610,4.296) -- (9.890,4.296);
\draw[lad tick,draw=ladMuted] (10.010,4.331) -- (10.290,4.331);
\draw[lad tick,draw=ladMuted] (7.840,5.326) -- (8.120,5.326);
\draw[lad lev,draw=ladInk2] (13.900,3.000) -- (15.100,3.000);
\draw[lad lev,draw=ladF] (13.900,3.800) -- (15.100,3.800);
\draw[lad pair,draw=ladInk2] (13.900,4.600) -- (15.100,4.600);
\draw[lad pair,draw=ladInk2] (13.900,1.200) -- (15.100,1.200);
\draw[lad lev,draw=ladI] (13.900,0.400) -- (15.100,0.400);
\draw[lad tick,draw=ladMuted] (14.290,6.160) -- (15.100,6.160);
\draw[lad tick,draw=ladMuted] (14.570,7.340) -- (15.100,7.340);
\draw[lad arr,lad on,draw=ladE] (2.800,3.000) -- node[pos=0.52,rotate=90,lad alab] {C2} (2.800,16.996);
\draw[lad arr,lad on,draw=ladE] (3.080,3.000) -- node[pos=0.40,rotate=90,lad alab] {G1} (3.080,12.915);
\draw[lad arr,lad opt,draw=ladE] (3.360,12.915) -- node[pos=0.78,rotate=90,lad alab] {SUPER} (3.360,3.000);
\draw[lad arr,lad on,draw=ladX] (3.080,12.915) -- node[pos=0.36,rotate=90,lad alab] {G2} (3.080,16.996);
\draw[lad arr,lad opt,draw=ladE] (3.360,12.915) -- node[pos=0.62,rotate=90,lad alab] {KSTEP} (3.360,16.996);
\draw[lad arr,lad on,draw=ladE] (3.620,3.000) -- (3.620,14.300);
\draw[lad arr,lad on,draw=ladE] (3.800,3.000) -- node[pos=0.55,rotate=90,lad alab] {C1a--c} (3.800,14.998);
\draw[lad arr,lad on,draw=ladE] (3.980,3.000) -- (3.980,15.750);
\draw[lad heat] (4.400,15.750) -- (4.400,12.915);
\draw[lad arr,lad on,draw=ladH] (2.700,12.915) -- (2.350,12.065);
\draw[lad arr,lad on,draw=ladX] (5.100,12.915) -- (5.100,12.200);
\draw[lad arr,lad on,draw=ladI] (5.250,16.996) -- node[pos=0.50,rotate=90,lad alab] {K15} (5.250,15.846);
\draw[lad arr,lad on,draw=ladI] (5.800,15.846) -- node[pos=0.45,rotate=90,lad alab] {R11p} (5.800,12.915);
\draw[lad arr,lad zero,draw=ladV] (6.120,12.200) -- node[pos=0.62,rotate=90,lad alab] {PEN\_KR2} (6.120,15.846);
\draw[lad hv,lad on,draw=ladV] (5.880,12.200) -- node[pos=0.50,rotate=90,lad alab] {KR2\_RAD, 147 nm} (5.880,3.770);
\draw[lad heat] (5.880,3.770) -- (5.880,3.000);
\draw[lad arr,lad on,draw=ladV] (6.300,12.200) -- node[pos=0.42,sloped,lad alab] {K2F\_BR} (7.620,8.276);
\draw[lad arr,lad on,draw=ladV] (6.300,12.200) -- node[pos=0.42,sloped,lad alab] {NK7} (8.660,8.856);
\draw[lad arr,lad zero,draw=ladV] (6.300,12.200) -- node[pos=0.62,sloped,lad alab] {K2F\_F} (8.900,8.856);
\draw[lad arr,lad on,draw=ladX] (6.300,12.915) -- node[pos=0.40,sloped,lad alab] {G3, G3C} (9.140,8.856);
\draw[lad arr,lad on,draw=ladI] (10.450,13.246) -- node[pos=0.50,rotate=90,lad alab] {K108} (10.450,8.856);
\draw[lad arr,lad on,draw=ladI] (10.720,14.396) -- node[pos=0.62,rotate=90,lad alab] {K26} (10.720,8.856);
\draw[lad arr,lad on,draw=ladB] (8.520,8.856) -- (8.520,8.276);
\draw[lad hv,lad on,draw=ladB] (9.350,8.856) -- node[pos=0.50,rotate=90,lad alab] {G4, 248 nm} (9.350,3.856);
\draw[lad hv,lad on,draw=ladB] (9.750,8.856) -- node[pos=0.50,rotate=90,lad alab] {G4A, 272 nm} (9.750,4.296);
\draw[lad heat] (9.750,4.296) -- (9.750,3.800);
\draw[lad hv,lad on,draw=ladC] (10.150,8.841) -- node[pos=0.50,rotate=90,lad alab] {G4C, 275 nm} (10.150,4.331);
\draw[lad heat] (10.150,4.331) -- (10.150,3.800);
\draw[lad arr,lad on,draw=ladH] (10.930,8.841) -- (10.930,8.041);
\draw[lad hv,lad on,draw=ladT] (7.980,8.276) -- node[pos=0.64,rotate=90,lad alab] {C3, 420 nm} (7.980,5.326);
\draw[lad heat] (7.980,5.326) -- (7.980,3.800);
\draw[lad arr,lad on,draw=ladH] (7.520,8.276) -- (7.520,7.476);
\draw[lad arr,lad zero,draw=ladH] (8.400,8.276) -- (8.400,7.476);
\draw[lad arr,lad on,draw=ladE] (14.360,3.000) -- node[pos=0.62,rotate=90,lad alab] {DIS\_a3} (14.360,6.160);
\draw[lad arr,lad on,draw=ladE] (14.640,3.000) -- node[pos=0.66,rotate=90,lad alab] {DIS\_A1} (14.640,7.340);
\draw[lad heat] (15.100,7.340) -- (15.100,4.600);
\draw[lad arr,lad on,draw=ladE] (14.080,3.000) -- node[pos=0.50,rotate=90,lad alab] {ATT} (14.080,1.200);
\draw[lad arr,lad on,draw=ladF] (14.880,4.600) -- node[pos=0.25,rotate=90,lad alab] {FF} (14.880,3.000);
\node[lad axis,anchor=east] at (0.870,1.000) {\ensuremath{-}2};
\node[lad axis,anchor=east] at (0.870,3.000) {0};
\node[lad axis,anchor=east] at (0.870,5.000) {2};
\node[lad axis,anchor=east] at (0.870,7.000) {4};
\node[lad axis,anchor=east] at (0.870,9.000) {6};
\node[lad axis,anchor=east] at (0.870,11.000) {8};
\node[lad axis,anchor=east] at (0.870,13.000) {10};
\node[lad axis,anchor=east] at (0.870,15.000) {12};
\node[lad axis,anchor=east] at (0.870,17.000) {14};
\node[lad axis,anchor=east] at (0.870,18.000) {eV};
\node[lad head,anchor=center] at (4.475,18.000) {Kr and Kr\ensuremath{_2}};
\node[lad head,anchor=center] at (9.650,18.000) {KrF and Kr\ensuremath{_2}F};
\node[lad head,anchor=center] at (14.500,18.000) {F\ensuremath{_2}};
\node[lad lab,anchor=east] at (2.570,3.000) {Kr};
\node[lad lab,anchor=east] at (2.570,12.915) {Kr\ensuremath{^*} 9.915};
\node[lad lab,anchor=east] at (2.570,16.996) {Kr\ensuremath{^+} 13.996};
\node[lad lab2l,anchor=north west] at (5.030,12.150) {Kr\ensuremath{_2^*}\\9.2};
\node[lad lab2,anchor=south east] at (6.300,15.896) {Kr\ensuremath{_2^+}\\12.846};
\node[lad tlab,anchor=west] at (4.460,14.300) {11.30};
\node[lad tlab,anchor=west] at (4.460,14.998) {11.998};
\node[lad tlab,anchor=west] at (4.460,15.750) {12.75};
\node[lad alab,anchor=east] at (2.310,12.195) {G3X};
\node[lad alab,anchor=east] at (5.040,12.557) {AK8\_Kr};
\node[lad hlab,anchor=west] at (5.980,3.370) {0.77};
\node[lad lab,anchor=west] at (12.380,14.396) {Kr\ensuremath{^+} + F\ensuremath{^-} 11.396};
\node[lad lab,anchor=west] at (12.380,13.246) {Kr\ensuremath{_2^+} + F\ensuremath{^-} 10.246};
\node[lad lab,anchor=south east] at (12.300,3.840) {Kr + F 0.8};
\node[lad lab2l,anchor=west] at (11.130,8.856) {KrF(B) 5.856\\KrF(C) 120\,cm\ensuremath{^{-1}} below\\BC, CB mix them};
\node[lad lab2,anchor=east] at (7.220,8.276) {Kr\ensuremath{_2}F\ensuremath{^*}\\5.276};
\node[lad slab,anchor=east] at (8.470,8.566) {G5,\\G5C};
\node[lad hlab,anchor=north] at (9.350,3.420) {X, 0.06};
\node[lad hlab,anchor=north] at (9.750,3.720) {A, 0.50};
\node[lad hlab,anchor=north] at (10.150,3.420) {A, 0.53};
\node[lad alab,anchor=north west] at (10.620,7.991) {QF2, QKR, QE\\QF2C, QKRC, QEC};
\node[lad hlab,anchor=north] at (7.980,3.720) {1.53};
\node[lad alab,anchor=north] at (7.520,7.426) {Q2F};
\node[lad alab,anchor=north] at (8.400,7.426) {Q2FF};
\node[lad lab,anchor=east] at (13.820,3.000) {F\ensuremath{_2}};
\node[lad lab,anchor=east] at (13.820,3.800) {F 0.8};
\node[lad lab,anchor=east] at (13.820,4.600) {F + F 1.6};
\node[lad lab,anchor=east] at (13.820,1.200) {F\ensuremath{^-} + F \ensuremath{-}1.8};
\node[lad lab,anchor=east] at (13.820,0.400) {F\ensuremath{^-} \ensuremath{-}2.6};
\node[lad tlab,anchor=west] at (15.150,6.160) {{\unsansmath\ensuremath{a\,^3\Pi_u}} 3.16};
\node[lad tlab,anchor=west] at (15.150,7.340) {{\unsansmath\ensuremath{A\,^1\Pi_u}} 4.34};
\draw[lad arr,lad on,draw=ladInk2] (1.250,-0.550) -- (1.950,-0.550);
\node[lad lab,anchor=west] at (2.050,-0.550) {a row of the default network};
\draw[lad arr,lad opt,draw=ladInk2] (8.850,-0.550) -- (9.550,-0.550);
\node[lad lab,anchor=west] at (9.650,-0.550) {a row of an optional group, off by default};
\draw[lad arr,lad zero,draw=ladInk2] (1.250,-0.970) -- (1.950,-0.970);
\node[lad lab,anchor=west] at (2.050,-0.970) {a band whose default is 0, no measurement};
\draw[lad hv,lad on,draw=ladB] (8.850,-0.970) -- (9.550,-0.970);
\node[lad lab,anchor=west] at (9.650,-0.970) {a photon, with its row and band};
\draw[lad arr,lad on,draw=ladH] (1.250,-1.390) -- (1.950,-1.390);
\node[lad lab,anchor=west] at (2.050,-1.390) {a quench, the level's energy to heat};
\draw[lad heat] (8.850,-1.390) -- (9.550,-1.390);
\node[lad lab,anchor=west] at (9.650,-1.390) {heat as a lower state falls apart (eV) or cascades};
\draw[lad pair,draw=ladInk2] (1.250,-1.810) -- (1.950,-1.810);
\node[lad lab,anchor=west] at (2.050,-1.810) {a pair, the sum of two formation energies};
\draw[lad tick,draw=ladMuted] (8.850,-1.810) -- (9.550,-1.810);
\node[lad lab,anchor=west] at (9.650,-1.810) {a threshold, or where a photon leaves its lower state};
\end{tikzpicture}}

\makeatletter
\newcommand\wi[1]{\ifcsname wi@#1\endcsname\csname wi@#1\endcsname
  \else\errmessage{no v0_2 number '#1'. Rerun make_overview_v0_2.py}\fi}
\newcommand\wk[1]{\ifcsname wk@#1\endcsname\csname wk@#1\endcsname
  \else\errmessage{no v0_3 number '#1'. Rerun make_overview_v0_3.py}\fi}
\newcommand\lad[1]{\ifcsname lad@#1\endcsname\csname lad@#1\endcsname
  \else\errmessage{no ladder number '#1'. Rerun make_energy_ladder.py}\fi}
\makeatother

% chemistry and units
\newcommand{\unit}[1]{\ensuremath{\,\mathrm{#1}}}
\newcommand{\el}{\ensuremath{\mathrm{e}^-}}
\newcommand{\Krs}{Kr\ensuremath{^*}}
\newcommand{\Krp}{Kr\ensuremath{^+}}
\newcommand{\Krtp}{Kr\ensuremath{_2^+}}
\newcommand{\Krts}{Kr\ensuremath{_2^*}}
\newcommand{\KtF}{Kr\ensuremath{_2}F\ensuremath{^*}}
\newcommand{\Ft}{F\ensuremath{_2}}
\newcommand{\Fm}{F\ensuremath{^-}}
\newcommand{\ra}{\ensuremath{\to}\ }
\newcommand{\Edep}{\ensuremath{E_\mathrm{dep}}}
\newcommand{\gIP}{\ensuremath{g_0I_\mathrm{sat}/P_\mathrm{in}}}
\newcommand{\x}[1]{\ensuremath{\times}#1}
% a factor of the measurement as a percentage of it (0.89 -> 89), to the
% factor's own precision: \pctof{x} for two decimals, \pctof[1]{x} for three
\newcommand{\pctof}[2][0]{\pgfmathparse{100*#2}%
  \pgfmathprintnumber[fixed,precision=#1]{\pgfmathresult}}

\begin{document}
\sansmath

{\LARGE\bfseries How accurate LIGHT v0\_3 is}\par
\vspace{4pt}
{\large The basic network of LIGHT v0\_3, version \lad{version}, on 30
September 2026: v0\_2 with the kinetic EEDF in the loop, with its electron
side, its energy ladder and its checks}\par
\vspace{10pt}

\noindent\fbox{\parbox{\dimexpr\linewidth-2\fboxsep-2\fboxrule\relax}{%
\setlength{\parskip}{5pt}%
\textbf{Bottom line.} LIGHT v0\_3 reproduces the fitted KrF decays [1, 2]
within \wk{sc.train.lo} to \wk{sc.train.hi}\,\% rms and predicts decays
that no fit used [4] within \wk{sc.xu_figs5_7}\,\% rms, as v0\_2 does, and
at the Tahoma laser [5] it gives \pctof{\wk{tah.g0.x}}\,\% of the measured
small-signal gain of \wk{tah.g0.meas}\unit{\%/cm}, where v0\_2 gives
\pctof{\wi{tah.g0.x}}\,\%.

Its electrons, solved kinetically during the run, land within
\wk{el.abs.g0}\,\% of the gain and \wk{el.abs.te}\,\% of the electron
temperature of ML's self-consistent LIGHT-Boltzmann loop at Tahoma. It
gives \pctof{\wk{tah.isat.x}}\,\% of Tahoma's saturation intensity and
\pctof{\wk{tah.eta.x}}\,\% of its formation efficiency, and the C/B
emission ratio of Yu and Setser [3] is matched to \wk{sc.ys_ir300}\,\%.
No fit used Tahoma's numbers and no rate was set by a laser measurement.

Its largest misses are v0\_2's, on data that conflict with the fitted sets,
\wk{sc.hakuta_fig2} and \wk{sc.hakuta_fig3}\,\% rms on the C-A decays of
Hakuta et al.\ [6] and \wk{sc.bs_tau_b}\,\% on the KrF(B) lifetime of
Burnham and Searles [7], and the VUV yield of pure krypton at low pump,
\pctof{\wk{pure.vuv.x}}\,\% of Morozov's measurement [8].

{\footnotesize [1] Eden, Waynant, Searles and Burnham, Appl.\ Phys.\
Lett.\ 32, 733 (1978). [2] Gadomski, Xu, Setser and Nelson, Chem.\ Phys.\
Lett.\ 189, 153 (1992). [3] Yu and Setser, J.\ Phys.\ Chem.\ 94, 2934
(1990). [4] Xu, Gadomski and Setser, J.\ Chem.\ Phys.\ 99, 2591 (1993).
[5] Spectra Technology, \emph{Experimental Support for KrF Laser
Modeling}, final report to Los Alamos, 19 July 1989, OSTI 10136715.
[6] Hakuta et al., Jpn.\ J.\ Appl.\ Phys.\ 29, 1194 (1990). [7] Burnham and Searles, J.\ Chem.\ Phys.\ 67, 5967 (1977).
[8] Morozov, Heindl, Kr\"ucken, Ulrich and Wieser, J.\ Appl.\ Phys.\ 103,
103301 (2008).\par}}}
\medskip

LIGHT v0\_3 is v0\_2 with the kinetic EEDF in the loop. Its chemistry,
its bands and its input are v0\_2's, and the v0\_2 overview describes
them. What changes is where the electron rates come from. v0\_2 integrates
every electron rate over a Maxwellian at $T_e$ with a tail read from a
table. v0\_3 solves the two-term Boltzmann equation with
electron-electron collisions during the run, at LIGHT's own mean energy,
and takes every electron rate, the elastic loss and the energy a captured
electron carries from that solution. The \wk{sc.same} kinetic datasets
have no electrons, so v0\_3 scores them exactly as v0\_2, and
\texttt{make\_overview\_v0\_3.py} refuses to write a number that differs.
The tests with electrons are Tahoma's, the pure-krypton yield and the
electron closure, and Table~\ref{tab:score} lists every test.

\begin{table}[h]
\centering\footnotesize
\renewcommand{\arraystretch}{1.12}
\begin{tabularx}{\linewidth}{@{}>{\raggedright\arraybackslash}X
  >{\raggedright\arraybackslash}p{2.3cm}>{\raggedright\arraybackslash}p{2.3cm}
  >{\raggedright\arraybackslash}p{2.3cm}>{\raggedright\arraybackslash}p{1.9cm}@{}}
\toprule
test & measured & v0\_3 & v0\_2 & role of the data\\
\midrule
energy ledger closure & & below \wk{closure.max} of \Edep{} in
  \wk{closure.runs} runs & below \wi{closure.max} & bookkeeping\\
KrF(B) decay against Kr and F$_2$, Eden 1978 & & rms \wk{sc.eden_kr} and
  \wk{sc.eden_f2}\,\% & the same & train\\
KrF(B,C) slow decay, Gadomski 1992 Figs.\ 1 and 2 & & rms
  \wk{sc.ksu_fig1} and \wk{sc.ksu_fig2}\,\% & the same & train\\
C/B emission ratio at 300\unit{K}, Yu and Setser 1990 & \wi{ys.meas}
  $\pm$ \wi{ys.err} & \wk{sc.ys_ir300}\,\% & the same & train\\
KrF(B,C) decays, Xu 1993 Figs.\ 5 to 7 & & rms \wk{sc.xu_figs5_7}\,\% &
  the same & held out\\
decay constants Xu 1993 printed, 50 and 100 Torr Kr & \wi{xu.50.meas}
  and \wi{xu.100.meas}\unit{s^{-1}} & \x{\wk{xu.50}} and \x{\wk{xu.100}} &
  the same & structure check\\
C-A decays, Hakuta 1990 & & rms \wk{sc.hakuta_fig2} and
  \wk{sc.hakuta_fig3}\,\% & the same & held out, in conflict\\
$\tau_B$, Burnham and Searles 1977 & \wi{sc.bs.meas} $\pm$
  \wi{sc.bs.err}\unit{ns} & \wk{sc.bs_tau_b}\,\% & the same & held out, in
  conflict\\
C/B ratio 230\unit{K} over 300\unit{K}, Yu and Setser & \wi{ysr.meas}
  $\pm$ \wi{ysr.err} & \wk{sc.ys_ratio_230_300}\,\% & the same & held out\\
\midrule
VUV yield of pure Kr at low pump, Morozov 2008 & \wk{pure.vuv.meas}
  $\pm$ \wk{pure.vuv.err} & \wk{pure.vuv}, \x{\wk{pure.vuv.x}} &
  \wi{pure.vuv}, \x{\wi{pure.vuv.x}} & check\\
electron closure at Tahoma, ML's kinetic loop & $T_e$
  \wk{el.loop.te}\unit{eV} & \wk{el.own.te}\unit{eV} &
  \wi{el.light.te}\unit{eV} & reference\\
\midrule
Tahoma small-signal gain $g_0$ & \wk{tah.g0.meas}\unit{\%/cm} &
  \wk{tah.g0}, \x{\wk{tah.g0.x}} & \wi{tah.g0}, \x{\wi{tah.g0.x}} &
  device, score only\\
\quad the same, Kannari's two \Krts{} estimates & & \wk{tah.kannari.g0},
  \x{\wk{tah.kannari.x}} & \wi{tah.kannari.g0}, \x{\wi{tah.kannari.x}} &
  device, score only\\
\quad the same, \Krts{} + F at the top of its band & & \wk{tah.hiF.g0},
  \x{\wk{tah.hiF.x}} & \wi{tah.hiF.g0}, \x{\wi{tah.hiF.x}} & device,
  score only\\
\quad the same, trimer on $T_g^{-1/2}$ & & \wk{tah.tri.g0},
  \x{\wk{tah.tri.x}} & \wi{tah.tri.g0}, \x{\wi{tah.tri.x}} & device,
  score only\\
Tahoma saturation intensity $I_\mathrm{sat}$ & \wk{tah.isat.meas} $\pm$
  \wi{tah.isat.err}\unit{MW/cm^2} & \wk{tah.isat}, \x{\wk{tah.isat.x}} &
  \wi{tah.isat}, \x{\wi{tah.isat.x}} & device, score only\\
Tahoma formation efficiency \gIP & \wk{tah.eta.meas} & \wk{tah.eta},
  \x{\wk{tah.eta.x}} & \wi{tah.eta}, \x{\wi{tah.eta.x}} & device, score
  only\\
\bottomrule
\end{tabularx}
\caption{LIGHT v0\_3 against the data, beside v0\_2, every band at its
default unless a row says otherwise. A percentage is the deviation from
the data, signed for a single value and rms for a set. A \x{} factor is
the model over the measurement. ``The same'' is checked: the generator
stops if v0\_3's number differs from v0\_2's.}
\label{tab:score}
\end{table}

\section{What v0\_3 changes}

The kinetic EEDF comes from \texttt{v0\_3/boltzmann.py}, ML's two-term
scheme with its electron-electron term, written on sparse matrices. It
carries exactly the electron processes the basic network books, so its
energy balance and LIGHT's ledger close the same budget. Penning
ionization injects its electrons at the energy the network gives them,
\Krtp{} recombination removes them with $\sigma v \propto \varepsilon^n$,
the template of its $T_e^n$ law, and the growth term carries the net. A
beam's secondary electrons are injected with the sub-excitation spectrum
of the degradation cascade, as ML's beam-sustained solver injects them.
The elastic and field powers LIGHT reads are the parts of the solver's own
Scharfetter-Gummel fluxes, so they are the energy it exchanges.

Every solution has LIGHT's mean energy, $\tfrac32 T_e$, held by a field
$E/N\ge0$, which in the pulse stands for the power the electrons take and
after it for the energy they still carry. Where even the solution without
a field is hotter than LIGHT's electrons, the field is 0. LIGHT reads a
table of the kinetic rates, captured energies and elastic power at five
mean energies around its own, renewed when any input of the kinetic state
drifts, and past the table's ends it carries them with its own Maxwellian
dependence on $T_e$. In the pulse the EEDF is advanced in time, one
backward-Euler step from the last update's to LIGHT's mean energy, with
the population's own change in the growth term. After the pulse it is
solved steady, and there the growth term is dropped where it cannot give
LIGHT's mean energy, where it moves more energy than the heating brings
and the gas does not relax the shape first, and where it has no positive
solution. \texttt{v0\_3/README.md} gives the design and what was found on
the way.

\section{How the numbers were made}

Every number in this note is written by
\texttt{docs/make\_overview\_v0\_3.py} from v0\_3 and from \texttt{ML/}, and
a test fails when those files no longer match the model. ML runs v0\_3
when \texttt{LIGHT\_VERSION} names it. As for v0\_2, Tahoma and pure
krypton run under the prototype beam drive, a test fixture in
\texttt{docs/} and not part of the model, whose ion pairs, excited atoms
and secondary-electron energy enter LIGHT's equations. Its secondary
electrons now enter the kinetic equation as well. ML's Tahoma
configuration is v0\_2's, the default network with every band at its
default and $n_e$ free from a seed of \wi{tah.seed}\unit{cm^{-3}}, now
with the kinetic EEDF in the loop. Each run is a process of its own, and
a Tahoma run takes minutes where v0\_2 takes seconds.

\section{The electron side at Tahoma}

ML's kinetic loop runs LIGHT on its Maxwellian closure, scales each
electron row by the ratio of a beam-sustained Boltzmann solution's rate
to the Maxwellian's at the same mean energy, and iterates until the
factors stop changing. v0\_3 does the same inside the run, at every
moment. Over Tahoma's \wi{tah.win.lo} to \wi{tah.win.hi}\unit{ns} window
the two agree:

\begin{table}[H]
\centering\footnotesize
\begin{tabular}{@{}lrrrr@{}}
\toprule
 & $T_e$ (eV) & $n_e$ (cm$^{-3}$) & $g_0$ (\%/cm) & formation\\
\midrule
ML's loop, \wk{el.passes} passes & \wk{el.loop.te} & \wk{el.loop.ne} &
  \wk{el.loop.g0} & \wk{el.loop.form}\\
v0\_3 & \wk{el.own.te} & \wk{el.own.ne} & \wk{el.own.g0} &
  \wk{el.own.form}\\
v0\_3 against the loop, \% & \wk{el.dev.te} & \wk{el.dev.ne} &
  \wk{el.dev.g0} & \wk{el.dev.formation}\\
v0\_2 against the loop, \% & \wi{el.dev.te} & \wi{el.dev.ne} &
  \wi{el.dev.g0} & \wi{el.dev.formation}\\
\bottomrule
\end{tabular}
\caption{v0\_3's electrons against ML's self-consistent loop at Tahoma.
v0\_2's deviations are absolute, v0\_3's signed.}
\end{table}

In the window LIGHT's mean energy is the kinetic one at \wk{tah.reached}
update, held by a field of at most \wk{tah.en}\unit{Td}, since the beam's
secondaries alone heat the electrons almost to it. The two energy
balances agree to \wk{tah.gap} of the beam's electron power.
Table~\ref{tab:factors} shows what the kinetic EEDF does to each electron
row, as its rate over LIGHT's Maxwellian rate at the same $T_e$.

\begin{table}[H]
\centering\footnotesize
\begin{tabular}{@{}lrr@{}}
\toprule
row & reference point, \wk{fac.ref.t}\unit{ns}, $T_e$ \wk{fac.ref.te}\unit{eV}
  & Tahoma, \wk{fac.tah.t}\unit{ns}, $T_e$ \wk{fac.tah.te}\unit{eV}\\
\midrule
\wkFactors
\bottomrule
\end{tabular}
\caption{The kinetic rate over the Maxwellian rate at the same $T_e$, for
each electron row the kinetic EEDF carries: G1 and C1a to C1c excite Kr,
C2 ionizes it, ATT attaches to \Ft{}, DIS\_a3 and DIS\_A1 dissociate it,
VIB1 and VIB4 are its first and last vibrational rows and R11p is \Krtp{}
recombination.}
\label{tab:factors}
\end{table}

\section{Tahoma and the bands}

Under the beam drive v0\_3 gives $g_0=$ \wk{tah.g0}\unit{\%/cm}, the
saturation intensity \wk{tah.isat}\unit{MW/cm^2} and \gIP{} $=$
\wk{tah.eta}, against v0\_2's \wi{tah.g0}, \wi{tah.isat} and
\wi{tah.eta}. Each band alone at its two ends moves the window's gain as
Table~\ref{tab:tahbands} shows. Of the bands, \wk{tah.bd.widest} spans the
most, as in v0\_2.

\begin{table}[H]
\centering\footnotesize
\begin{tabular}{@{}lll@{}}
\toprule
band & $g_0$ (\%/cm), low to high end & over the measured $g_0$\\
\midrule
\wkTahBands
\bottomrule
\end{tabular}
\caption{Each band alone at its two ends at Tahoma under v0\_3, the
others at their defaults, over the \wi{tah.win.lo} to
\wi{tah.win.hi}\unit{ns} window. With every band at its default these
runs give \wk{tah.win.g0}\unit{\%/cm}. ML's formation run, the same
network with ML's formation channels added, gives the \wk{tah.g0} of
Table~\ref{tab:score}.}
\label{tab:tahbands}
\end{table}

\section{The reference point}

At the notes' reference point, \wk{ref.p}\unit{atm},
\wk{ref.xf2}\,\% \Ft{}, \wk{ref.P}\unit{kW\,cm^{-3}} for
\wk{ref.pulse}\unit{\mu s} and run to \wk{ref.tend}\unit{\mu s}, v0\_3
and v0\_2 give these shares of the deposited energy:

\begin{table}[H]
\centering\footnotesize
\begin{tabular}{@{}lrrr@{}}
\toprule
 & v0\_3 & v0\_2 & v0\_3 over v0\_2, \%\\
\midrule
248\unit{nm}, KrF B-X & \wk{ref.248} & \wk{ref.248.v02} & \wk{ref.248.d}\\
272\unit{nm}, KrF B-A & \wk{ref.272} & \wk{ref.272.v02} & \wk{ref.272.d}\\
275\unit{nm}, KrF C-A & \wk{ref.275} & \wk{ref.275.v02} & \wk{ref.275.d}\\
420\unit{nm}, \KtF{} & \wk{ref.420} & \wk{ref.420.v02} & \wk{ref.420.d}\\
147\unit{nm}, \Krts{} & \wk{ref.147} & \wk{ref.147.v02} & \wk{ref.147.d}\\
all bands & \wk{ref.all} & \wk{ref.all.v02} & \\
to \Krs{}, gross & \wk{ref.krstar} & \wk{ref.krstar.v02} & \\
to \Krp{} & \wk{ref.krplus} & \wk{ref.krplus.v02} & \\
KrF(B) formed & \wk{ref.krfb} & \wk{ref.krfb.v02} & \\
elastic heating & \wk{ref.elastic} & \wk{ref.elastic.v02} & \\
$T_g$ at the end (K) & \wk{ref.tg} & \wk{ref.tg.v02} & \\
\bottomrule
\end{tabular}
\caption{The reference point, percent of \Edep{} unless a row says
otherwise. v0\_2's numbers are v0\_3's run on v0\_2's closure, which is
v0\_2 (a test checks it).}
\end{table}

The kinetic electrons spend less on elastic heating and ionization and
more on exciting Kr, and more KrF(B) forms. Over the pulse's second half
the field brings \wk{ref.ratio.lo} to \wk{ref.ratio.hi} times the input
power. The excess is the table's renewal. Between renewals LIGHT runs on
a table made at the last renewal's electron density and gas. So at each
renewal its electrons sit off the fresh table's balance by the excess,
and they regain that balance well before the next renewal. Renewed every
\wk{conv.dt.ns}\unit{ns}, the excess is at most \wk{ref.ratio.dt.hi}.
With LIGHT's own $n_e\,d\langle\varepsilon\rangle/dt$ at the renewal, less
the step's own, the two balances agree to \wk{ref.gap} of the input
power. The run takes \wk{ref.solutions} kinetic solutions, with $E/N$
between \wk{ref.en.lo} and \wk{ref.en.hi}\unit{Td}.
Renewing the table at half the drift moves no band by more than
\wk{conv.step} of itself, renewing it every \wk{conv.dt.ns}\unit{ns} by
\wk{conv.dt}, and doubling the energy cells by \wk{conv.cells}.

\section{The first nanosecond}
\label{sec:first}

In the reference pulse's first \wk{first.t}\unit{ns} the electrons heat
from the gas temperature, and most of the input power goes into raising
their mean energy, which no steady solution holds. Solved steady, as
v0\_3 solved it until 1 October 2026, the EEDF dropped the growth term at
\wk{first.steady.goff} of its \wk{first.steady.n} updates there, its
field carried \wk{first.steady.ratio.lo} to \wk{first.steady.ratio.hi} of
the input power, and its energy balance missed LIGHT's by up to
\wk{first.steady.gap} of it. In the pulse v0\_3 now advances the EEDF in
time, one backward-Euler step from the last update's to LIGHT's mean
energy, with the population's own change in the growth term and no rule
to drop it. In the first nanosecond the field then carries
\wk{first.mem.ratio.lo} to \wk{first.mem.ratio.hi} of the input power,
\wk{first.mem.goff} of the \wk{first.n} updates drop the growth term, and
after \wk{first.t.split}\unit{ns} the two energy balances agree to
\wk{first.mem.gap}. Before that, while the electrons are thermal and a
cell is as wide as $kT$, they agree to \wk{first.mem.gap0}.

The bands barely notice. The steady first nanosecond moves no band at
the reference point by more than \wk{first.d.steady} of itself, and
LIGHT's own Maxwellian closure in it moves none by more than
\wk{first.d.maxwell}, since only \wk{first.edep} of the pulse's energy
enters then. At Tahoma a steady start moves the window's gain by
\wk{first.tah.d} and $n_e$ by \wk{first.tah.ne.d}, and
\wk{first.tah.steady.unreached} of its first nanosecond's steady
solutions missed LIGHT's mean energy, against
\wk{first.tah.mem.unreached} of the time steps.

\section{Pure krypton}

Under the beam drive pure krypton at Morozov's low pump gives a VUV yield
of \wk{pure.vuv} with v0\_3, \x{\wk{pure.vuv.x}} the measured
\wk{pure.vuv.meas}, against v0\_2's \wi{pure.vuv}. At
\wk{pure.hi.P}\unit{kW\,cm^{-3}} it gives \wk{pure.hi}, against v0\_2's
\wi{pure.hi.def}. The kinetic
electrons do not close v0\_2's gap at low pump, which stays either physics
LIGHT lacks or a systematic error of the measurement.

\section{The kinetic EEDF's own checks}

\begin{itemize}
\item With ML's electron-electron quadrature the solver is ML's dense one
(\texttt{ee.solve\_field}) to \wk{ml.dF} in $F$ and \wk{ml.dk} in every
rate at \wk{ml.states} states with the growth term. v0\_3's own quadrature
moves $F$ by \wk{ml.own} at most there.
\item The growth frequency is the solution's own net frequency, and the
energy a field and the Penning electrons bring equals what the
collisions, the captures and the growth term take, to \wk{bal.pulse} in
the pulse and \wk{bal.nofield} with no field. Without the growth term the
net loss of electrons comes back in the first cell, and with that
injection the balance closes to \wk{bal.nogrowth}, a residual first order
in the cell width.
\item With no drive the solution is the grid's Maxwellian at the gas
temperature to \wk{und.600.dev}, and it exchanges no energy with the gas
to round-off. Taken by central differences, the elastic power gave it a
loss to the gas, and late in an afterglow LIGHT's electrons cooled below
the kinetic solution's mean. With ML's quadrature, $F$ constant in
each cell, the undriven solution condenses into the first cell, to a mean
energy of
\wk{und.600.cell}\unit{eV} on \wk{und.lo} cells and
\wk{und.2400.cell}\unit{eV} on \wk{und.hi}, where the gas gives
\wk{und.32kt}\unit{eV}.
\item \wk{s.t}\unit{ns} into the reference pulse, at
\wk{s.field}\unit{Td}, the problem without the growth term has two
solutions, \wk{s.hot}\unit{eV} from a hot start and \wk{s.cold}\unit{eV}
from a cold one, and at a fixed mean energy it has one. The field that
mean energy needs falls from \wk{s.en.006}\unit{Td} at
\wk{s.m.006}\unit{eV} to \wk{s.en.02}\unit{Td} at \wk{s.m.02}\unit{eV}
and rises to \wk{s.en.05}\unit{Td} at \wk{s.m.05}\unit{eV}.
\item The beam's secondaries reach the grid with the injected mean energy
to \wk{inj.dev} of it.
\item A time step too long to matter is the steady growth solution to
\wk{step.big}, and one too short to matter, with no field, leaves $F$
where it was to \wk{step.tiny}. With e-e left out, a picosecond's step
from the gas's Maxwellian under a Joule power of
\wk{step.power}\unit{eV/s} per electron raises the mean energy by what
the field and the sources bring less what the collisions, the captures
and the population's change take, to \wk{step.energy}. With e-e at
thermal energies, where a cell is as wide as $kT$, the flux scheme's
energy exchange holds only to \wk{step.ee600} on \wk{und.lo} cells and
\wk{step.ee2400} on \wk{und.hi}.
\end{itemize}

\section{The energy ladder}

v0\_3's network is v0\_2's, and so is its ladder (Figure~\ref{fig:ladder}).
KrF(B) lies \lad{DeB}\unit{eV} below the ion pair \Krp{} + \Fm{} and
\lad{BoverKrF}\unit{eV} above Kr + F, and the \lad{wl.G4}-nm photon
carries \lad{hv.G4}\unit{eV} of that.

\begin{figure}[p]
\centering
\ladpicture
\caption{The energy ladder of LIGHT v0\_3, version \lad{version}, in eV
above ground-state Kr and \Ft{} at rest. A bar is a species at its
formation energy and a dashed bar a pair, the sum of two. An arrow is a
row, labelled with its id. It is solid in the default network and dotted
for a band whose default is 0, \lad{zero}. The rows of the optional group
\lad{groups} are dashed. A photon's arrow drops by its quantum, and the red
squiggle below it is the heat its lower state gives the gas as it falls
apart. A red stub is a quench. The energy a row releases heats the gas,
except a photon's, which leaves, and that of \lad{toe}, which goes to the
electrons. KrF(B) and KrF(C) lie \lad{gapBC}\unit{cm^{-1}} apart and are
drawn as a pair of bars that every B and C row meets. The rows
\lad{notdrawn} move nothing between levels and are not drawn.}
\label{fig:ladder}
\end{figure}

\section{Where v0\_3 is not yet reliable}

\begin{itemize}
\item Everything in v0\_2's list, since the chemistry is v0\_2's: \Krts{}
+ F, the trimer above \wi{tri.klimek.hi}\unit{K}, \Krs{} + 2Kr above its
measured range, hot gas, ion-ion recombination near atmospheric pressure,
the two Penning rates, \Krts's well, and Hakuta's and Burnham and Searles'
conflicts.
\item The steady shape after the pulse. There the kinetic EEDF is a
steady solution, while the electrons' shape is set by their history, and
v0\_3 takes the steady shape at their mean energy. In the pulse it now
follows the history (Section~\ref{sec:first}).
\item The first cell. Where the growth term is dropped, after the pulse
only, the net loss of electrons comes back cold in the first cell, which
cools a solution at low field. v0\_3 holds LIGHT's mean energy regardless,
so this moves the field and the low-energy shape.
\item Thermal electrons. On the default \wk{und.lo} cells the grid's
Maxwellian at \wk{und.tg}\unit{K} has a mean energy of
\wk{und.600.mean}\unit{eV}, below the \wk{und.32kt}\unit{eV} of the gas,
so late in an afterglow LIGHT's electrons sit a few percent below the gas
temperature. \wk{und.hi} cells give \wk{und.2400.mean}\unit{eV}.
\item Cost. A run takes minutes, so the bands' scans and the scan over the
notes' ranges are hours under the kinetic closure.
\end{itemize}

\appendix
\section{Reaction rates}

v0\_3's rows, rates and bands are v0\_2's, and the v0\_2 overview's
Appendix~A lists them. Its electron-impact rows run at the kinetic EEDF's
rates, which Table~\ref{tab:factors} compares with the Maxwellian ones.

\end{document}
